\documentclass[aspectratio=169, 11pt]{beamer}
\usepackage[utf8]{inputenc}
\usepackage[T1]{fontenc}
\usepackage[french]{babel}
\usepackage{graphicx}
\usepackage{tikz}
\usepackage{listings}
\usepackage{xcolor}
\usepackage{hyperref}

% Thème propre et moderne
\usetheme{Madrid}
\usecolortheme{beaver} % Un peu moins "bleu standard", plus pro

% Header simple : titre en noir, sans bandeau rouge
\setbeamercolor{frametitle}{fg=black,bg=white}
\setbeamercolor{title in head/foot}{fg=black,bg=white}

% Footer : projet + noms (gauche) et numéro de page (droite)
\setbeamertemplate{footline}{%
  \leavevmode%
  \hbox{%
    \begin{beamercolorbox}[wd=.82\paperwidth,ht=2.25ex,dp=1ex,leftskip=1.2ex]{author in head/foot}%
      \usebeamerfont{author in head/foot}%
      \footnotesize RVP2 — Moteur de recherche multimodal (Ping) — Barbier / Ledoux / Pascal / Wone
    \end{beamercolorbox}%
    \begin{beamercolorbox}[wd=.18\paperwidth,ht=2.25ex,dp=1ex,right,rightskip=1.2ex]{page number in head/foot}%
      \usebeamerfont{page number in head/foot}%
      \footnotesize \insertframenumber/\inserttotalframenumber
    \end{beamercolorbox}%
  }%
}

% Raccourci pour une "pause" (révélation progressive)
\newcommand{\pauseSlide}{\pause}

% Configuration des blocs de code
\lstset{
    basicstyle=\scriptsize\ttfamily,
    keywordstyle=\color{blue!80!black}\bfseries,
    stringstyle=\color{red!60!black},
    commentstyle=\color{gray},
    frame=single,
    breaklines=true,
    backgroundcolor=\color{gray!5},
    showstringspaces=false
}

\title[Architecture Technique]{Deep Dive Technique : Backend & IA}
\subtitle{Embeddings, Indexation et Moteur de Recherche}
\author{Nathan Barbr - Projet Informatique}
\date{\today}

\setbeamertemplate{headline}{%
  \leavevmode%
  \begin{beamercolorbox}[wd=\paperwidth,ht=2.5ex,dp=1.125ex,leftskip=1.2ex,rightskip=1.2ex]{section in head/foot}%
    \usebeamerfont{section in head/foot}%
    \insertsectionhead%
    \ifx\insertsubsectionhead\empty\relax\else%
      \hfill{\footnotesize\insertsubsectionhead}%
    \fi%
  \end{beamercolorbox}%
}

\begin{document}

% =============================================================================
% PAGE DE GARDE (style Centrale Lyon)
% =============================================================================
\begin{frame}[plain]
    \vspace*{-0.2cm}
    \begin{minipage}[t]{0.55\textwidth}
        \includegraphics[height=1.2cm]{Logo_Centrale_Lyon.png}
    \end{minipage}%
    \begin{minipage}[t]{0.44\textwidth}
        \raggedleft
    \end{minipage}

    \vfill

    \begin{center}
        {\LARGE\bfseries RVP2 : Moteur de recherche\\[0.25em]
        multimodal de vidéos de\\[0.25em]
        tennis de table\par}
    \end{center}

    \vfill

    \begin{flushright}
        \small
        BARBIER Nathan\\
        LEDOUX Edgar\\
        PASCAL Romain\\
        WONE Mamoudou
    \end{flushright}
\end{frame}

% =============================================================================
% SOMMAIRE (logo en haut à droite uniquement ici)
% =============================================================================
\begin{frame}[plain]{Sommaire}
    \begin{tikzpicture}[remember picture,overlay]
        \node[anchor=north east, xshift=-0.4cm, yshift=-0.2cm] at (current page.north east)
        {\includegraphics[height=1.0cm]{Logo_Centrale_Lyon.png}};
    \end{tikzpicture}

    % Mise en page plus compacte et mieux équilibrée (évite les gros retraits)
    \vspace{0.3cm}
    {\small
    \setlength{\leftmargini}{1.2em}
    \setlength{\leftmarginii}{1.6em}
    \setlength{\itemsep}{0.35em}
    \setlength{\parsep}{0pt}
    \setlength{\topsep}{0.2em}
    \setlength{\partopsep}{0pt}
    \setlength{\parskip}{0pt}
    \begin{itemize}
        \setlength\itemsep{0.8em}
        \item \textbf{Contexte \& objectifs}
        \item \textbf{Démonstration}
        \item \textbf{Aspects techniques}
              \begin{itemize}
                  \setlength\itemsep{0.3em}
                  \item Embeddings \& Pipeline d’indexation
                  \item Recherche Hybride (Elasticsearch)
                  \item Focus IA \& Perspectives
              \end{itemize}
        \item \textbf{Gestion de projet}
              \begin{itemize}
                  \setlength\itemsep{0.3em}
                  \item Organisation \& Planning
                  \item Méthodologie (Agile, GitFlow, PR)
              \end{itemize}
        \item \textbf{Conclusion}
    \end{itemize}
    }
\end{frame}

% =============================================================================
% SECTION : GESTION DE PROJET (GANTT)
% =============================================================================
% =============================================================================
% SLIDE : CONTEXTE (après le sommaire)
% =============================================================================
\section{Contexte \& objectifs}
\begin{frame}{Contexte : Expertise LIRIS en Tennis de Table}
    \begin{block}{Collaboration JO Paris 2024}
        \begin{itemize}\large
            \item Le laboratoire LIRIS : expertise reconnue
            \item Participation aux JO de Paris 2024 avec l'équipe de France
            \item Doctorant en appui de l'analyse vidéo
        \end{itemize}
        
        \vspace{0.8cm}
        \begin{center}
            \large\textbf{Comment rendre accessible au grand public des outils d'analyse vidéo\\
            professionnels utilisés en compétition olympique ?}
        \end{center}
    \end{block}
\end{frame}

% =============================================================================
% SLIDE : OBJECTIFS (après contexte)
% =============================================================================
\begin{frame}{Objectifs}
    \begin{block}{Objectif}
        \begin{itemize}\large
            \item Avoir une plateforme web pour faciliter les recherches dans les matchs de ping-pong,
            pouvant intégrer différentes modalités d’analyses (sonore, visuelle, texte, métadonnées...).
        \end{itemize}
    \end{block}

    \vspace{0.4cm}
    \begin{block}{Différents personas}
        \begin{itemize}\large
            \item Entraîneurs
            \item Joueurs
            \item Fans
        \end{itemize}
    \end{block}
\end{frame}

% =============================================================================
% SLIDE : DONNÉES PRÉSENTES
% =============================================================================
\begin{frame}{Données présentes}
    \begin{columns}[T,onlytextwidth]
        \begin{column}{0.55\textwidth}
            \centering
            % Image fournie par l'utilisateur (à ajouter au projet sous le nom image.png)
            \includegraphics[width=\linewidth]{these.png}
        \end{column}
        \begin{column}{0.45\textwidth}
            \vspace{0.2cm}
            \begin{block}{Source}
                \begin{itemize}\large
                    \item Documents et données issus de la thèse d’Aymeric ERADES
                \end{itemize}
            \end{block}
            \vspace{0.3cm}
            \begin{block}{Contenu des données}
                \begin{itemize}\large
                    \item Match découpé en points
                    \item Service utilisé
                    \item Type de coups, etc.
                \end{itemize}
            \end{block}
        \end{column}
    \end{columns}
\end{frame}

% =============================================================================
% SLIDE : INTRO DÉMO
% =============================================================================
\section{Démonstration}
\begin{frame}{Démo : Recherche dans un match}
    \begin{center}
        \LARGE\textbf{Maintenant : démonstration du projet}\\[0.6cm]
        \Large\href{http://localhost:3000}{\texttt{http://localhost:3000}}
    \end{center}
\end{frame}

% =============================================================================
% SLIDE : DONNÉES PRÉSENTES
% =============================================================================
\section{Coulisses techniques}


% =============================================================================
% SECTION 2 : LES EMBEDDINGS & L'INDEXATION
% =============================================================================
\subsection{Embeddings \& Pipeline d'indexation}

\begin{frame}{Vue d'Ensemble du Flux de Données}
    \centering
    \includegraphics[width=\linewidth,height=0.55\textheight,keepaspectratio]{tikz.png}

    \vspace{0.1cm}

    \begin{block}{Points Clés}
        \begin{itemize}\small
            \item \textbf{Flux unidirectionnel} pour l'indexation (Data $\to$ IA $\to$ Base).
            \item \textbf{Elasticsearch} agit comme source de vérité unique (Texte + Vecteurs).
            \item \textbf{FastAPI} sert d'interface d'interrogation unifiée.
        \end{itemize}
    \end{block}
\end{frame}

\begin{frame}{1. Structure d'un Document Indexé}
    \textbf{Ce que nous stockons réellement dans Elasticsearch} :
    Une fusion de 3 types de données pour chaque point joué.
    
    \vspace{0.2cm}
    
    \begin{columns}
        \begin{column}{0.3\textwidth}
            \begin{block}{A. Métadonnées (CSV)}
            \textit{Données déterministes}\\
            \begin{itemize}
                \item \texttt{match\_id} (Keyword)
                \item \texttt{set\_num} (Int)
                \item \texttt{winner} (Keyword)
                \item \texttt{nb\_coups} (Int)
                \item \texttt{faute\_type} (Keyword)
            \end{itemize}
            \end{block}
        \end{column}
        \begin{column}{0.35\textwidth}
            \begin{block}{B. Texte Généré}
            \textit{Pour la recherche floue}\\
            Champ \texttt{description} :\\
            "Point du set 3. Echange de 12 coups remporté par Fan Zhendong sur un topspin revers."
            \end{block}
        \end{column}
        \begin{column}{0.3\textwidth}
            \begin{block}{C. Vecteur (IA)}
            \textit{Pour le sens}\\
            Champ \texttt{embedding} :\\
            \texttt{[0.12, -0.05, 0.88, ...]}
            \\(384 dimensions)
            \end{block}
        \end{column}
    \end{columns}
\end{frame}

\begin{frame}[fragile]{2. Pipeline d'Indexation (ETL)}
    \textbf{Fichier :} \texttt{backend/embeddings/indexer.py}
    
    Le processus est séquentiel et déterministe :
    
    \begin{enumerate}
        \item \textbf{Extraction} : Lecture du CSV ligne par ligne.
        \item \textbf{Enrichissement} : 
        \begin{itemize}
            \item Transformation des codes (`g1`, `m2`) en texte naturel ("gauche courte", "milieu mi-table").
            \item Génération de la phrase descriptive complète.
        \end{itemize}
        \item \textbf{Vectorisation} : Appel du modèle \texttt{Sentence-Transformer} sur la description.
        \item \textbf{Indexation} : Envoi du document JSON complet (Méta + Texte + Vecteur) vers Elasticsearch via l'API \texttt{\_bulk}.
    \end{enumerate}
\end{frame}

% =============================================================================
% SECTION 2 : LES EMBEDDINGS & L'INDEXATION
% =============================================================================


% =============================================================================
% SECTION 3 : LOGIQUE DE RECHERCHE
% =============================================================================
\subsection{Recherche Hybride}

\begin{frame}{Avantages d'Elasticsearch contre Apache Solr}
    \begin{itemize}
        \setlength\itemsep{1em}
        \item Recherche textuelle
        \item Highlights
        \item Autocomplétion
        \item Format des données input en JSON schema-less
        \item Requêtes de type body
    \end{itemize}
    
    \vspace{1cm}
    \centerline{\large $\to$ Plus adapté à nos besoins}
\end{frame}

\begin{frame}{Étapes d'une requête utilisateur (recherche textuelle)}
    \textit{Exemple : ``Topspin revers set 3''}
    
    \vspace{0.3cm}
    \begin{enumerate}
        \item \textbf{Connexion} : le backend ouvre un socket vers le container Docker ES (port 9200).
        \item \textbf{Filtrage} : ES élimine d'abord tous les documents qui ne sont pas du \texttt{set 3}.
        \item \textbf{Matching} : sur les documents restants, ES cherche \texttt{topspin} \textbf{ET} \texttt{revers}.
        \item \textbf{Scoring (BM25)} : ES note les résultats (un point avec ``topspin puissant revers'' score mieux qu’un simple ``revers'').
        \item \textbf{Agrégations} : en parallèle, ES calcule les stats (combien de points gagnants dans ce sous-ensemble ?) pour mettre à jour les facettes de l’UI.
    \end{enumerate}
\end{frame}

\begin{frame}{Scénario 1 : Recherche Textuelle \& Filtres}
    \textit{L'utilisateur cherche : "Revers gagnant" dans le Set 3}
    
    \textbf{Processus Backend (`routers/search.py`) :}
    
    \begin{enumerate}
        \item \textbf{Réception} : L'API reçoit \texttt{q="Revers gagnant", set\_num=3}.
        \item \textbf{Construction de la Query (Booléenne)} :
            \begin{itemize}
                \item \textbf{MUST} (Score) : \texttt{multi\_match} sur les champs texte avec "Revers gagnant".
                \item \textbf{FILTER} (Binaire) : \texttt{term} exact sur \texttt{set\_num : 3}.
            \end{itemize}
        \item \textbf{Exécution} : Elasticsearch filtre d'abord (très rapide), puis score les documents restants.
    \end{enumerate}
\end{frame}

\begin{frame}{Scénario 2 : La Feature "Highlights" (Recherche Vectorielle)}
    \textit{L'utilisateur clique sur le bouton "Highlights" \textrightarrow{} Comment ça marche ?}
    
    \textbf{Concept : Recherche par "Vecteur Moyen Idéal" (`routers/semantic.py`)}
    
    L'algorithme ne cherche pas juste "des points longs", il cherche des points qui \textit{ressemblent} aux meilleurs moments.
    
    \begin{enumerate}
        \item \textbf{Définition du "Gold Standard"} (Interne) : 
        Le système identifie d'abord les points "parfaits" existants (ex: Points gagnants avec $>8$ coups).
        \item \textbf{Calcul du Profil} : 
        On calcule la \textbf{moyenne vectorielle} de ces 50 meilleurs points. Cela crée un "Vecteur Highlight" théorique.
        \item \textbf{Recherche k-NN} :
        On utilise ce vecteur moyen comme requête pour trouver les points les plus proches sémantiquement dans l'espace vectoriel.
        
        $\Rightarrow$ \textit{Résultat : On trouve des points intenses (ex: 7 coups, sauvetage incroyable) même s'ils ne respectent pas strictement les règles CSV.}
    \end{enumerate}
\end{frame}

\begin{frame}{Pourquoi cette architecture ?}
    \begin{itemize}
        \item \textbf{Précision} : Les filtres exacts (Set, Joueur) garantissent qu'on ne se trompe pas de match.
        \item \textbf{Découverte} : La recherche vectorielle ("Highlights", "Similar") permet de trouver des points basés sur le \textit{style} de jeu plutôt que des mots-clés exacts.
        \item \textbf{Performance} : Tout est indexé au même endroit (Elasticsearch). Pas de latence réseau entre une base SQL et une base Vectorielle.
    \end{itemize}
\end{frame}

\subsection{Focus IA \& Perspectives}

\begin{frame}{Pourquoi introduire de l'IA ?}
    \begin{columns}[T, onlytextwidth]
        \begin{column}{0.48\textwidth}
            \begin{block}{Limites actuelles (SQL)}
                \begin{itemize}
                    \setlength\itemsep{1em}
                    \item \textbf{Trop rigide} : Mots-clés exacts uniquement.
                    \item \textbf{Aveugle} : Ignore synonymes et contexte.
                    \item \textbf{Binaire} : Pas de notion d'intensité ou de style.
                \end{itemize}
            \end{block}
        \end{column}
        \begin{column}{0.48\textwidth}
            \begin{block}{Apport de l'Intelligence Artificielle}
                \begin{itemize}
                    \setlength\itemsep{1em}
                    \item \textbf{Compréhension} : Interprète l'intention utilisateur.
                    \item \textbf{Sémantique} : Trouve le sens au-delà des mots.
                    \item \textbf{Similarité} : Recommande des points visuellement proches.
                \end{itemize}
            \end{block}
        \end{column}
    \end{columns}
\end{frame}

\begin{frame}{Etat Avancement}
    \begin{columns}[T, onlytextwidth]
        \begin{column}{0.55\textwidth}
            \textbf{\large Compréhension des requêtes (LLM)}
            \vspace{0.3cm}
            \begin{block}{Objectif}
                Transformer une phrase libre en structure exploitable.
            \end{block}
            \vspace{0.3cm}
            \textbf{Exemple :}
            \begin{center}
                \textit{"Montre-moi les points gagnants en topspin de Hugo"}
                
                $\downarrow$
                
                \texttt{JSON structuré}
            \end{center}
        \end{column}
        \begin{column}{0.40\textwidth}
            \textbf{\large Implémentation}
            \vspace{0.3cm}
            \begin{itemize}
                \setlength\itemsep{0.5em}
                \item Ollama
                \item \texttt{mistral:7b-instruct}
                \item Structured output JSON
                \item Validation et fallback
            \end{itemize}
            \vspace{0.6cm}
            \textbf{\textrightarrow{} Le LLM agit comme un parser intelligent.}
        \end{column}
    \end{columns}
\end{frame}

\begin{frame}{Limites \& Perspectives IA avancées}
    \textbf{\Large Limites actuelles}
    \vspace{0.6cm}
    \begin{itemize}
        \setlength\itemsep{1.2em}
        \item Pas encore de recherche hybride fine (bool + vector pondéré)
        \item Pas de re-ranking LLM
        \item Embedding uniquement textuel
        \item Pas d'explicabilité
    \end{itemize}
\end{frame}

\begin{frame}{Limites \& Perspectives IA avancées}
    \textbf{\Large Perspectives (1/2)}
    \vspace{0.4cm}
    
    \begin{block}{1. Hybrid Search}
        Combinaison : 
        \vspace{0.2cm}
        \centerline{$ Score\_final = \alpha \cdot score\_bool + \beta \cdot score\_vec $}
        \vspace{0.2cm}
        \textrightarrow{} Meilleur contrôle et précision.
    \end{block}
    
    \vspace{0.6cm}
    
    \begin{block}{2. Re-ranking LLM}
        \begin{itemize}
            \item kNN $\to$ Top 20
            \item LLM $\to$ Re-classement intelligent
        \end{itemize}
    \end{block}
\end{frame}

\begin{frame}{Limites \& Perspectives IA avancées}
    \textbf{\Large Perspectives (2/2)}
    \vspace{0.4cm}
    
    \begin{block}{3. Recherche par Embedding Textuel Avancé}
        Enrichissement des descriptions :
        \begin{itemize}
            \item Résumé automatique
            \item Intensité du point
            \item Contexte match
            \item Style de jeu
        \end{itemize}
        \vspace{0.2cm}
        \textrightarrow{} \textit{Espace sémantique plus riche.}
    \end{block}
    
    \vspace{0.6cm}
    
    \textbf{4. Recherche par Embedding Vidéo (Multimodal)}
\end{frame}
% =============================================================================
% SECTION : GESTION DE PROJET (GANTT)
% =============================================================================
\section{Gestion de projet}

\begin{frame}{Répartition des tâches}
    \begin{block}{Responsabilités Techniques}
        \begin{itemize}
            \setlength\itemsep{1em}
            \item \textbf{Embeddings} : Génération et optimisation des vecteurs.
            \item \textbf{Elasticsearch} : Configuration du cluster et indexation.
            \item \textbf{Backend} : API, gestion des routes et logique métier.
            \item \textbf{Frontend} : Interface utilisateur et expérience (UX).
            \item \textbf{LLM / IA} : Intégration Ollama, Hugging Face et pipelines RAG.
        \end{itemize}
    \end{block}
\end{frame}

\begin{frame}{Diagramme de Gantt}
    \begin{columns}[T,onlytextwidth]
        \begin{column}{0.38\textwidth}
            \begin{block}{Objectif}
                \small Visualiser l’avancement et les dépendances entre lots.
            \end{block}
            \vspace{0.15cm}
            \begin{block}{Organisation}
                \begin{itemize}\small
                    \item \textbf{Sprint 1--2} : données, pipeline embeddings, indexation.
                    \item \textbf{Sprint 3--4} : recherche hybride + API.
                    \item \textbf{Sprint 5} : UX, perf, déploiement, démo.
                \end{itemize}
            \end{block}

        \end{column}
        \begin{column}{0.62\textwidth}
            \centering
            % Plus de place à l'image (lisibilité)
            \includegraphics[width=\linewidth,height=0.84\textheight,keepaspectratio]{diagramme_gantt.png}
        \end{column}
    \end{columns}
\end{frame}

\begin{frame}{Agile}
    \begin{columns}[T,onlytextwidth]
        \begin{column}{0.62\textwidth}
            \centering
            \includegraphics[width=\linewidth,height=0.78\textheight,keepaspectratio]{agile.png}
        \end{column}
        \begin{column}{0.38\textwidth}
            \begin{block}{Notre organisation}
                \begin{itemize}\small
                    \item Cycles courts : planifier \textrightarrow{} développer \textrightarrow{} tester.
                    \item Revue régulière : validation avec le besoin utilisateur.
                    \item Déploiement itératif : chaque sprint ajoute une fonctionnalité.
                \end{itemize}
            \end{block}
        \end{column}
    \end{columns}
\end{frame}

\begin{frame}{GitFlow (workflow Git)}
    \begin{columns}[T,onlytextwidth]
        \begin{column}{0.55\textwidth}
            \begin{block}{Branches}
                \begin{itemize}\small
                    \item \textbf{\texttt{main}} : versions stables / démos.
                    \item \textbf{\texttt{develop}} : intégration continue des features.
                    \item \textbf{\texttt{feature/*}} : développement d’une fonctionnalité (PR \textrightarrow{} \texttt{develop}).
                    \item \textbf{\texttt{hotfix/*}} : correctifs urgents (merge vers \texttt{main} et \texttt{develop}).
                \end{itemize}
            \end{block}
            \vspace{0.2cm}
            \begin{block}{Règles}
                \begin{itemize}\small
                    \item Commits petits + messages clairs.
                    \item Revue de code avant merge.
                    \item Tag des versions lors des jalons.
                \end{itemize}
            \end{block}
        \end{column}
        \hfill
        \begin{column}{0.40\textwidth}
            \begin{block}{Cycle typique}
                \begin{enumerate}\small
                    \item Créer \texttt{feature/x}
                    \item Dev + tests
                    \item PR \textrightarrow{} \texttt{develop}
                    \item Release \textrightarrow{} \texttt{main}
                \end{enumerate}
            \end{block}
        \end{column}
    \end{columns}
\end{frame}

\begin{frame}{Schéma GitFlow}
    \centering
    \includegraphics[width=\linewidth,height=0.85\textheight,keepaspectratio]{gitflow.png}
\end{frame}

% =============================================================================
% CONCLUSION
% =============================================================================
\begin{frame}{Processus de Pull Request (PR)}
    \begin{columns}[T,onlytextwidth]
        \begin{column}{0.48\textwidth}
            \begin{block}{Workflow de validation}
                \begin{itemize}\small
                    \setlength\itemsep{1em}
                    \item \textbf{Description} : Résumé clair des changements (Feature/Fix).
                    \item \textbf{Checklist} : Lint, Tests, Doc.
                    \item \textbf{Revue} : Validation par au moins un pair avant merge.
                \end{itemize}
            \end{block}
            \vspace{0.5cm}
            \small \textbf{\textrightarrow{} Garantit la qualité du code sur \texttt{develop}.}
        \end{column}
        \hfill
        \begin{column}{0.48\textwidth}
            \centering
            \vspace{-0.2cm}
            \includegraphics[width=\linewidth,height=0.75\textheight,keepaspectratio]{pr.png}
        \end{column}
    \end{columns}
\end{frame}

% =============================================================================
% CONCLUSION
% =============================================================================
\section{Conclusion}

\begin{frame}{Bilan du projet}
    \begin{columns}[T,onlytextwidth]
        \begin{column}{0.48\textwidth}
            \begin{block}{Accomplissements}
                \begin{itemize}\small
                    \item \textbf{Innovation} : Recherche hybride (Texte + Vecteurs) fonctionnelle.
                    \item \textbf{Performance} : Architecture scalable avec Elasticsearch.
                    \item \textbf{UX} : Interface fluide et adaptée aux experts du ping.
                \end{itemize}
            \end{block}
        \end{column}
        \begin{column}{0.48\textwidth}
            \begin{block}{Perspectives : Vers une IA Expert}
                \begin{itemize}\small
                    \item \textbf{Langage Naturel} : Requêter la base de matchs avec des phrases complexes et fluides.

                    \item \textbf{Déploiement} : Ouverture de la plateforme à la communauté.
                \end{itemize}
            \end{block}
        \end{column}
    \end{columns}
\end{frame}

\begin{frame}[plain]
    \centering
    \vspace{1.5cm}
    
    {\Huge \textbf{Merci de votre attention !}}
    
    \vspace{1cm}
    
    {\Large Avez-vous des questions ?}
    
    \vspace{2cm}
    
    \small
    \textit{Projet informatique RVP2 — 2025-2026}
\end{frame}

\end{document}
